\documentclass[UTF8,11pt]{ctexart}
\usepackage[a4paper,margin=22mm]{geometry}
\usepackage{amsmath,amssymb,bm,graphicx,adjustbox,fvextra,longtable,array}
\xeCJKDeclareCharClass{CJK}{"2160 -> "217F, "2460 -> "249B, "2200 -> "22FF}
\newcommand{\fitmath}[1]{\begin{adjustbox}{max width=\linewidth}$\displaystyle #1$\end{adjustbox}}
\setlength{\parindent}{0pt}\setlength{\parskip}{.65em}\renewcommand{\arraystretch}{1.3}
\sloppy\allowdisplaybreaks
\begin{document}
\section*{2016年考研政治答案解析}
\subsection*{2016 年思想政治理论试题答案及解析}

\subsection*{一、单项选择题}

（每小题1分，共16分）


\subsection*{1.【标准答案】B}

本题考查适度原则。“度”是指保持事物质的稳定性的数量界限，即事物的范围、幅度和限度，这一原理要求我们坚持适度的原则，防止“过犹不及”。题干中“加一点点盐”导致“菜美味极了”，而“一把盐”放进嘴里，“结果又苦又咸”。从“一点盐”到“一把盐”所导致的不同后果，说明“盐”必须适度，不能少也不能多，故B正确。A干扰性最大，注意：题干并非是因为往“同一个菜里”持续加盐而导致“菜美味极了”变为“又苦又咸”，故 A不符合题意。C表述错误，我们能通过现象去认识事物的本质。D表述虽然正确，但题干并未体现这一点。事物的自我否定，是指事物的内部矛盾运动。


\subsection*{2.【标准答案】A}

本题考查马克思主义的自由观。自由是标示人的活动状态的范畴，是人在活动中通过认识和利用必然所表现出的一种自觉自主的状态。自由是有前提条件的：\textcircled{1}不以牺牲别人的自由为条件；\textcircled{2}坚持“必然”（自然规律）为前提条件。如果夸大必然否定自由将导致“宿命论”，夸大自由否定必然将导致“唯意志论”。题干第一种观点既主张自由又不否认自然规律为前提，属于唯物辩证法的观点，第二种观点过分夸大自由的作用，属于唯意志论的观点。故 A正确。


\subsection*{3.【标准答案】D}

本题考查剩余价值率。这是一道计算题，可采取“设问→条件”的方式解答。首先，剩余价值率用公式表示为：\(\displaystyle m'=\dfrac{\displaystyle m}{\displaystyle v}\)，因此，只需要求出 m（剩余价值）和v（可变资本）即可。其次，根据题干可知：\textcircled{1}资本家投资100万元，其中，购买劳动力支付10万元，也即 \(\displaystyle v=10\)万元；\textcircled{2}一轮生产后总资本为120元，这意味着该资本家获得了 \(\displaystyle 120-100=20\)万元的利润，因为剩余价值与利润在数量上是相等的，这就意味着 \(\displaystyle m=20\)万元。综合两个条件，最后可以计算出 \(\displaystyle m'=\dfrac{\displaystyle 20}{\displaystyle 10}=200\%\)，故D正确。


\subsection*{4.【标准答案】A}

本题考查金融垄断资本的发展。20世纪 70年代初，由于资本主义发展不平衡的加深和国际货币体系内在矛盾的激化，布雷顿森林体系崩溃，随后西方国家普遍走上了金融自由化和金融创新的道路。金融自由化与金融创新是金融垄断资本得以形成和壮大的重要制度条件，推动着资本主义经济的金融化程度不断提高，故A正确。BCD 虽然对当代西方资本主义金融资本快速发展壮大也有一定程度的影响，但并非“重要制度条件”，与题意不符。


\subsection*{5.【标准答案】C}

本题考查马克思主义中国化理论成果的精髓。实事求是既是党的思想路线的核心，也是马克思主义两大理论成果的精髓，贯穿了马克思主义两大理论成果始终，故C正确。ABD不符合题意。A“改革创新”是时代精神的核心。B“独立自主”D“群众路线”和C“实事求是”一起构成毛泽东思想活的灵魂。


\subsection*{6.【标准答案】B}

本题考查《论十大关系》的内容。1956年，毛泽东发表了《论十大关系》，从十个方面论述了我国社会主义建设需要重点把握的重大关系，并指出：“这十个问题都是围绕一个基本方针，就是把国内外一切积极因素调动起来，为社会主义服务事业服务。”故B正确。A是1957年发表的《关于正确处理人民内部矛盾的问题》的内容；C是1956 年中共八大确定的经济建设方针；D是1956年1月召开的知识分子问题会议上提出来的。


\subsection*{7.【标准答案】B}

本题考查社会主义初级阶段理论。这是一道阅读理解题，题干包含两层含义：\textcircled{1}“新世纪以来，我国经济和社会发展呈现新特征”，这说明社会主义初级阶段发展的阶段性特征；\textcircled{2}“这些特征并未改变我国处于社会主义初级阶段的基本事实”，这说明我国处于并将长期处于社会主义初级阶段，故B正确。AD表述正确，但题干并未体现。C讲的是“社会主义初级阶段的矛盾”，与“社会主义初级阶段”是两个不同问题。


\subsection*{8.【标准答案】D}

本题考查民族区域自治的核心。民族区域自治是在统一而不可分割的国家领导下，在各少数民族聚居的地方设立自治机关，行使自治权，实行区域自治。其核心是保障少数民族当家做主，管理本民族、本地方事务权利，故D正确。其他三项均不符合题意。A是党的统一战线政策的重要内容。B是处理民族问题的基本原则。C是党的民族政策的重要内容。


\subsection*{9.【标准答案】B}

本题考查民族意识的觉醒。近代以来，帝国主义的侵略给中华民族带来了巨大灾难，也从反面教育了中国人民。鸦片战争以后，先进的中国人开始睁眼看世界。中日甲午战争后，中国人民的民族意识开始普遍觉醒，故B正确。ACD 均不符合题意。A带来的结果是全民族抗战意识兴起；C使中国人民认识到清王朝已经沦为洋人的朝廷；D是中国近代史上第一次不败而败的战争。


\subsection*{10.【标准答案】C}

本题考查救亡图存。魏源在《海国图志》中提出了“师夷长技以制夷”的思想，这里的“长技”是指“西方的军事和科学技术”，即通过学习外国的先进军事和科学技术，以期富国强兵，抵御外国侵略，故C正确。ABD均不符合题意。


\subsection*{11.【标准答案】D}

本题考查抗日战争中的重要战役。忻口会战是抗日战争初期，在华北战场上规模最大、最激烈的一次战役，也是国共两党军队合作抗日、配合最好的一次战役。故D正确。A发生在1938年3月，是抗战爆发后中国正面战场取得的首次重大胜利。B发生在全面抗战打响之前（1933年）。C发生在1948年秋，是解放战争时期三大战役之一。


\subsection*{12.【标准答案】C}

本题考查解放战争时期的土地政策。在解放战争时期，《五四指示》将党在抗日战争时期实行的减租减息政策变为“耕者有其田”政策，这标志着解放区在农民土地问题上，开始由抗日战争时期的削弱封建剥削，向变革封建土地关系、废除封建剥削过渡，故C正确。A属于新中国成立初期在解放区进行土改的政策。B是王明“左”倾教条主义的表现。D是1928年《井冈山土地法》的政策。


\subsection*{13.【标准答案】B}

本题考查爱国主义的特点。爱国主义是历史的、具体的，在不同的历史条件和文化背景下所形成的爱国主义，总是具有不同的内涵和特点。故B正确。爱国主义是“客观的”、“抽象的”、“现实的”三种表述均是错误的。


\subsection*{14.【标准答案】A}

本题考查权利保障。新大纲将“人权保障”的表述变更为“权利保障”。权利保障主要是指对公民权利的法律保障，具体包括公民权利的宪法保障、立法保障、行政保护和司法保障。其中，宪法保障是权利保障的前提和基础；立法保障是权利保障的重要条件；行政保护是权利保障的关键环节，司法保障是公民权利保障的最后防线。故 A正确。


\subsection*{15.【标准答案】C}

\subsection*{16.【标准答案】B}

\subsection*{二、多项选择题 （每小题2分，共 34分。多选或少选均不得分。）}

\subsection*{17.【标准答案】ACD}

本题考查唯物论、辩证法、认识论相关的原理。这是一道“遍地开花”题，先要抓住核心主旨——-显微镜下的“一沙一世界”，然后从哲学各个角度进行分析。\textcircled{1}从唯物论角度看，世界统一于物质，但世界的统一性是多样性的统一，世界上的一切事物都具有多样化的属性，故 A正确。\textcircled{2}从辩证法的角度看，人的认识是由个别上升到一般，再由一般到个别的辩证发展过程，题干中“一沙一世界”即是从个别（一粒沙）到一般（多彩世界）的认识过程，故D正确。\textcircled{3}从认识论角度看，实践为认识发展提供了必要的条件，通过日益完备的物质手段强化了主体的认识能力，题干中的显微摄影就是人们制造出来并用来认识世界的工具，这说明通过制造和使用工具，人们可以加深对客观世界的认识，故C正确。B表述错误，本质是事物的根本性质，是相对稳定的，不会随着人们认识的变化而变化。


\subsection*{18.【标准答案】ABD}

本题考查时间的一维性。时间的一维性，是指时间永远按照过去、现在、将来一个方向演变，一去不返。题干中的诗句表现出物是人非的历史变迁感，体现了时间的一维性。A“物换星移几度秋”感叹时光的不经意流逝，B“莫待无花空折枝”提醒人们时不我待、光阴易逝，D“黑发不知勤学早，白首方悔读书迟”感叹岁月易逝、光阴一去不复返，都具有时间一维性特点，故 ABD正确。C的大意是“云起日落风吹”，说明雨要来了，这是对景观的描述，不涉及时间一维性。


\subsection*{19.【标准答案】BCD}

本题考查资本主义经济危机理论。具体有：\textcircled{1}生产过剩是资本主义经济危机的本质特征。这种过剩是相对过剩，是指相对于劳动人民有支付能力的需求来说社会生产的商品显得过剩（群众想消费，但消费不起），而不是与劳动人民的实际需求相比的绝对过剩（群众不想消费），故 A错误，D正确；\textcircled{2}资本主义的基本矛盾是资本主义经济危机的根本原因，这一矛盾表现在：生产无限扩大的趋势与劳动人民有支付能力的需求相对缩小的矛盾；个别企业内部生产的有组织性和整个社会生产的无政府状态之间的矛盾，故BC 正确。


\subsection*{20.【标准答案】ABD}

本题考查经济全球化的动因。主要有：\textcircled{1}科学技术的进步和生产力的发展为经济全球化提供了坚实的物质基础和根本的推动力；\textcircled{2}跨国公司的发展为经济全球化提供了适宜的企业组织形式；\textcircled{3}各国经济体制变革是经济全球化的体制保障。故 ABD 正确。C是个别企业自身行为，对推动经济全球化迅猛发展影响甚微。


\subsection*{21.【标准答案】CD}

本题考查社会主义从空想到科学。空想社会主义的局限性主要表现在：\textcircled{1}他们只看到了资本主义必然灭亡的命运，却未能揭示资本主义必然灭亡的经济根源；\textcircled{2}他们要求埋葬资本主义，却看不到埋葬资本主义的力量；\textcircled{3}他们憧憬取代资本主义的理想社会，却找不到通往理想社会的现实道路。科学社会主义之所以能够超越空想社会主义，原因就在于马克思恩格斯在新的历史条件下创立了唯物史观和剩余价值学说，揭示了资本主义必然灭亡的经济根源，找到了实现理想社会的现实道路，故CD正确。A是科学社会主义与空想社会主义的共同点，不符合题意。B表述错误，科学社会主义没有对未来社会进行细致的描述。


\subsection*{22.【标准答案】ABC}

本题考查全面提高对外开放水平。中国高铁走出国门、走向世界意味我国融入世界经济的程度在加深，国际投资合作水平日益提升，参与国际竞争能力明显增强，自主创新能力显著提高，故 ABC 正确。我国传统产业的调整还在进行中，未得到根本性改变，D错误。


\subsection*{23.【标准答案】AC}

本题考查社会主义初级阶段的基本经济制度。这是一道阅读理解题，题目中“保障民企平等使用土地、减免税收”的政策有利于保证各种所有制经济依法平等使用生产要素，“扩大民间投资在电网、电信、铁路等非竞争性领域的参与力度”的政策有利于保证各种所有制经济公开、公平、公正参与市场竞争，故 AC 正确。BD是企业内部管理体制和资产结构的问题，与题意无关。


\subsection*{24.【标准答案】ABCD}

本题考查精准扶贫的措施。十八届五中全会指出：实施精准扶贫、精准脱贫，因人因地施策，提高扶贫实效。对于有劳动能力的，支持发展特色产业和转移就业；对于“一方水土养不起一方人”的，实施扶贫搬迁；对于生态特别重要和脆弱的，实行生态保护扶贫；对于丧失劳动能力的，实施兜底性保障政策。故 ABCD正确。


\subsection*{25.【标准答案】ABCD}

本题考查简政放权的目的。处理好政府与市场的关系，就是要使市场在资源配置中起决定性作用和更好地发挥政府的作用。此轮行政审批制度的改革是简政放权的重要体现，简政放权旨在减少审批环节，降低市场交易成本，激发市场主体的内在活力和社会创造力；让政府摆正自身位置，承担好服务市场的角色，加快完善社会主义市场经济体制，提高治理能力和治理水平。故 ABCD 正确。


\subsection*{26.【标准答案】ABC}

本题考查统一战线的意义。习近平指出：统一战线作为中国共产党不断夺取革命、建设和改革事业胜利的重要法宝，作为实现祖国统一和中华民族伟大复兴的重要法宝，决不能丢掉；作为党的一个政治优势，决不能削弱：作为党的一项长期方针，决不能动摇，故 ABC 正确。中国共产党的领导是人民当家作主的根本保证，D表述错误。


\subsection*{27.【标准答案】BCD}

本题考查维新变法失败的原因。主要有：\textcircled{1}客观上是由于以慈禧太后为首的强大的守旧势力的反对；\textcircled{2}主观上是由于民族资产阶级力量弱小和维新派自身的局限性。其局限性表现为：不敢否定封建主义；对帝国主义抱有幻想；惧怕人民群众，故 BCD 正确。A 表述错误，维新派仍要打着孔子的旗号，借古代圣贤之名推行“托古改制”，没有彻底否定“中学”。


\subsection*{28.【标准答案】ABC}

本题考查新文化运动的历史意义。新文化运动是中国历史上一次前所未有的启蒙运动和空前深刻的思想解放运动。其意义有：\textcircled{1}唤醒了一代青年，使中国的知识分子尤其是广大青年受到西方民主和科学思想的洗礼；\textcircled{2}掀起了一股生气勃勃的思想解放潮流，冲决了禁锢人们思想的闸门，激励着人们去探求救国救民的真理；\textcircled{3}涌现了一批青年革命民主主义者，接受俄国十月革命的影响，为马克思主义在中国的广泛传播准备了思想和文化条件，故 ABC 正确。新文化运动对孔学的批判是充满理性的，并没有彻底否定孔学的历史作用，D表述错误。


\subsection*{29.【标准答案】ABCD}

本题考查邓小平南方谈话的内容。主要包括：\textcircled{1}提出计划和市场都是经济手段。邓小平指出，计划多一点还是市场多一点，不是社会主义与资本主义的本质区别。计划不等于社会主义，资本主义也有计划；市场不等于资本主义，社会主义也有市场；\textcircled{2}提出社会主义的本质是解放生产力，发展生产力，消灭剥削，消除两极分化，最终达到共同富裕；\textcircled{3}提出判断改革开放的“三个有利于”标准。即：是否有利于发展社会主义社会的生产力；是否有利于增强社会主义国家的综合国力；是否有利于提高人民的生活水平；\textcircled{4}提出“发展才是硬道理”的重要论断；\textcircled{5}强调加强党的建设；\textcircled{6}提出社会主义经历一个长过程发展后必然代替资本主义。因此，ABCD 均正确。


\subsection*{30.【标准答案】ABCD}

本题考查法治体系的内涵。体现在四个方面：\textcircled{1}“法治体系”分为法律规范体系、法治实施体系、法治监督体系、法治保障体系和党内法规体系，即“法律体系”是“法治体系”的重要组成部分；\textcircled{2}“法律体系”重在法律及制度，要求有法可依、有法必依、执法必严、违法必究，而“法治体系”要实现“法治”，还要科学立法、严格执法、公正司法、全民守法；\textcircled{3}建设社会主义法治体系意味着，不仅要制定法律，更要实施法律、遵守法律、落实法律，为全面推进依法治国明确政治保证、体制机制、工作布局和努力方向；\textcircled{4}“法治体系”不仅仅是静态的法律文本，而且也是动态的法的实现过程。故 ABCD正确。


\subsection*{31.【标准答案】ABD}

本题考查法律权利和法律义务的关系。表现为：\textcircled{1}两者是相互依存的关系。法律权利的实现必须以相应法律义务的履行为条件，同样，法律义务的设定和履行也必须以法律权利的行使为根据；\textcircled{2}两者是目的与手段的关系。离开了法律权利，法律义务就失去了履行的价值和动力；离开了法律义务，法律权利也形同虚设；\textcircled{3}两者具有复合性的关系（注意：新教材已经把“二重性”改为“复合性”），即一个行为可以同时是权利行为和义务行为。故 ABD正确。由于权利和义务是相互依存的，不存在一方先出现，另一方再出现的情况，故C表述错误。


\subsection*{32.【标准答案】ABCD}

\subsection*{33.【标准答案】ABD}

\subsection*{三、分析题 （每小题10分，共50分）}

\subsection*{34.【标准答案】}

（1）实践是认识的基础。实践是认识的来源和认识发展的动力，是检验认识是否具有真理性的标准。“四个全面”来自实践，要发挥好它对实践的指导和引领作用，还须回到实践中去经受检验，以进一步发展完善。（2分）认识是主体在实践基础上对客体的能动反映，来自实践的认识或理论对实践具有指导作用。来自实践的四个全面”对新的实践又具有指导引领作用。（2分）“四个全面”的每一个“全面”都是实践——认识——再实践的一个动态、发展变化的、循环往复的演进过程。（1分）


（2）所谓辩证思维，就是指运用辩证唯物主义理论与方法，承认事物的联系与发展，承认矛盾、分析矛盾、解决矛盾，善于抓住关键、找准重点，洞察事物的发展规律。（2分）“四个全面”体现了事物普遍联系与发展的思想，揭示了“四个全面”之间及各个“全面”战略目标和举措之间的联系与发展。体现了矛盾分析方法，它既直指当前我国发展中面临的突出问题与主要矛盾，抓住了发展中的“牛鼻子”，又对各种矛盾心中有数，以通过优先解决主要矛盾和矛盾的主要方面，来带动其他矛盾的解决；还体现了整体地、全面地看待问题以及分析问题的方法，树立了全局观。（3分）（注：只要考生结合认识的本质（实践与认识的关系）、认识的辩证发展过程、辩证法的根本观点（联系的观点、发展的观点、全面的观点以及矛盾观点）、“四个全面”的内涵及其相互关系等内容作答，言之成理，均可酌情给分。）


\subsection*{35.【标准答案】}

（1）经济发展与环境治理之间应该是一种双赢而非对立的关系。在这种关系下，既要保证经济的持续健康发展，又要满足人民群众的基本环境需要，而不是单纯强调环境治理或经济发展的需要。（2分）一些地方、一些领域没有处理好经济发展同生态环境治理的关系，以无节制消耗能源、破坏环境为代价换取经济发展，导致能源资源、生态环境问题越来越突出，最终将会危及人民群众的基本环境需要。（2分）应正确处理经济发展同生态环境治理的关系，树立绿色发展理念，推动绿色发展。（2分）


（2）建设生态文明是一场涉及生产方式、生活方式、思维方式和价值观念的革命性变革。实现这样的变革，必须建立系统完整的制度体系，用制度保护生态环境，推进生态文明建设。（1分）把体现生态文明建设状况的指标纳人经济社会发展评价体系，使之成为推进生态文明建设的重要导向和约束；划定生态保护红线，完善生态补偿机制，建立责任追究制度；健全法律法规，完善生态环境保护管理制度。（3分）（注；经济发展和环境治理是提高人民生活水平的两个方面，经济要发展，环境也要保护，不是矛盾，我们要从生态文明制度建设上来处理好二者的关系。考生从生态文明建设以及绿色发展理念的角度回答问题，言之成理，均可酌情给分。）


\subsection*{36.【标准答案】}

（1）1949年春，夺取全国胜利已成定局，党的工作重心很快就要实现由乡村到城市战略转移，共产党即将从革命党成为执政党。（2分）新中国成立后，因为胜利，党面临着骄傲自满、“糖衣炮弹”攻击等危险；党还面临着恢复和发展经济，进行社会主义建设，实现国家富强人民幸福的艰巨任务。能否适应新形势、担负起新使命，是进城之后的中国共产党必须谨慎面对和要解决的问题，因此，1949年进北平实际上就是我党即将进行的另一场“考试”。（3分）


（2）进入新世纪以来，世情、国情、党情发生了深刻变化，党面临着长期执政、改革开放、市场经济、外部环境的考验，面临着精神懈怠、能力不足、脱离群众、消极腐败的危险，党的自身建设任务比过去任何时候都更为繁重、更为迫切。（2分）另一方面，党肩负着全面建成小康社会、实现中华民族伟大复兴的重任，肩负着不断坚持和发展中国特色社会主义的重任。党所肩负的重任与党所面临的新的、更大的考验与风险，要求党必须更加谦虚谨慎、戒骄戒躁、艰苦奋斗，所以说，“赶考”远未结束。（3分）（注：考生从七届二中全会召开的历史背景和意义及现阶段党的建设紧迫性和意义作答，言之有理，均可酌情给分。）


\subsection*{37.【标准答案】}

（1）家风通过家规、家教的方式得以传承并体现在家庭成员的日常行为中。好家风是中华民族优良文明基因的重要组成部分，是家庭或家族在代代相传中不断积淀而成的。（2分）好家风的传承既是塑造个体良好品德的过程，又是推动社会文明进步的过程。（2分）好家风与其他优秀文化传统的传承一样，是延续优良文明基因的过程。（1分）


（2）好家风与社会主义核心价值观是息息相通的，是涵养社会主义核心价值观的深厚土壤。材料中勤俭持家、友善邻里、廉洁自律的好家风正是社会主义核心价值观的具体体现。（2分）传承好家风，就是要通过延续其优良风尚，摒弃不合时代要求的内容，并进行创新性发展，使好家风成为提高个人素养、增进家庭和睦、培育良好社会风气的重要途径，让弘扬社会主义核心价值观真正落到实处。（3分）（注：考生围绕中华民族优良传统、家庭美德和社会主义核心价值观回答问题，言之成理，均可酌情给分。）


\subsection*{38.【标准答案】}

（1）联合国宪章宗旨和原则体现了世界人民渴望和平与平等、合作与繁荣的美好愿望；国际形势发生深刻变化，传统和非传统安全威胁相互交织；某些国家妄图挑战第二次世界大战以来的世界秩序和国际体系，使联合国面临前所未有的挑战。（3分）联合国是发展多边主义的最佳平台，是维护战后和平与稳定的基石；捍卫联合国的权威，强化联合国有效应对全球性威胁和挑战的能力，维护以联合国宪章宗旨和原则为核心的国际秩序和国际体系是实现国际社会长治久安、推动人类社会进步的根本。（2分）


（2）当今世界，各国相互依存，你中有我、我中有你，“打造人类命运共同体”意味着各国需要求同存异和休戚与共，以开放包容、合作共赢的心态谋求共同发展；以不断建设和完善机制性合作来发挥建设性作用，构建新型国际关系；（3分）这既是对联合国宪章倡导的主权平等、不干涉内政、和平解决国际争端等基本精神的继承与发展，符合和平与发展的时代主题和国际社会的共同利益需求，也体现出中国的世界眼光、博大胸襟和历史担当。（2分）

\end{document}
